\subsection{An explicit description of the septic residual field}

Let $A=\operatorname{Jac}(v^2+v=u^7)$,
$F=\mathbb Q(\zeta_7)$, $\lambda=1-\zeta_7$, and
$E:y^2+xy=x^3-x^2-2x-1$.
Retain the residual field and character calculations of the
preceding septic lemmas: $K=\mathbb Q(A[3])$ contains $F$,
$\operatorname{Gal}(K/F)=C_{56}$, its conductor divides
$(3)\lambda^2$, and the coefficient-norm identity controls
$E[3]$.

\begin{proposition}\label{prop:lm-explicit-septic-field}
There is an equality of subfields of an algebraic closure of
$\mathbb Q$,
\[
 K=\mathbb Q\bigl(\zeta_7,\sqrt[7]{3},E[3]\bigr).
\]
Its degree is $336$.  The intermediate field fixed by the
normal subgroup of order seven is
$F(E[3])$, of degree $48$ over $\mathbb Q$.
\end{proposition}
\begin{proof}
On the residual coefficient field $\mathbb F_{3^6}$, the
coefficient norm to $\mathbb F_{3^2}$ is $a\mapsto a^{91}$.
Polarization identifies the mod-three cyclotomic character
on $C_{56}$ with $a\mapsto a^{28}$.  The character identity
for $E$ therefore gives $a\mapsto a^{91-28}=a^7$ on
$C_{56}$.  Thus $F(E[3])/F$ is precisely its degree-eight
subextension.

The seven-primary part of the ray class group of $F$ for
the modulus $(3)\lambda^2$ has order at most seven.
Indeed, $F$ has class number one.  The seven-primary parts
of $(\mathcal O_F/(3))^\times$ and
$(\mathcal O_F/\lambda^2)^\times$ are each cyclic of
order seven.  The global group $\mu_7$ maps nontrivially
to both factors: reduction at three preserves order seven,
and $\zeta_7=1-\lambda$ is nontrivial modulo $\lambda^2$.
The quotient of these two local factors by its diagonal
image has order seven.  Further global units can only
make the ray quotient smaller.  The degree-seven
subextension of $K/F$ shows that it is not smaller.
There is consequently a unique cyclic extension of degree
seven with conductor dividing this modulus.

We verify directly that $F(\sqrt[7]{3})$ has this conductor.
It has degree seven, since the valuation of $3$ at the
unique prime over three is one.  It is cyclic because
$\mu_7\subset F$.  At three it is totally and tamely
ramified of degree seven, and hence has conductor exponent
one.  It is unramified at primes other than three and seven
by the Kummer equation.

At seven put $\alpha=\sqrt[7]{3}$ and $\beta=\alpha-3$.
The polynomial
\[
 (T+3)^7-3
\]
is Eisenstein over $\mathbb Q_7$: all its nonleading
coefficients are divisible by seven and
$v_7(3^7-3)=1$.  Thus
$L_7=\mathbb Q_7(\alpha)$ is totally ramified of degree
seven with $\mathcal O_{L_7}=\mathbb Z_7[\beta]$.
Its different has exponent seven, since the derivative
is $7\alpha^6$ and $\alpha$ is a unit.

Let $F_7=\mathbb Q_7(\zeta_7)$ and
$M_7=F_7L_7$.  The coprime totally ramified degrees six
and seven give $[M_7:\mathbb Q_7]=42$ and
$e(M_7/L_7)=6$.  This last extension is tame, with
different exponent five.  Transitivity of the different
therefore gives
\[
 d(M_7/\mathbb Q_7)=5+6\cdot7=47,
 \qquad
 d(M_7/F_7)=47-7\cdot5=12.
\]
For a cyclic degree-seven extension, all six nontrivial
characters have the same conductor and the
conductor--discriminant formula gives their common
conductor exponent $12/6=2$.
This proves the required local conductor at $\lambda$.

The uniqueness from the ray group now identifies
$F(\sqrt[7]{3})$ with the degree-seven subextension of
$K/F$.  Its degree is coprime to that of $F(E[3])/F$,
so their compositum is $K$.  The degree assertions
follow from $6\cdot7\cdot8=336$.
\end{proof}

\begin{corollary}[The split residual group]\label{cor:septic-residual-group}
The full residual Galois group is
\[
 \operatorname{Gal}(K/\Q)\simeq C_{56}\rtimes_3 C_6,
\]
where a generator of $C_6$ acts on $C_{56}$ by cubing.
In particular the extension by $\operatorname{Gal}(F/\Q)$ splits.
\end{corollary}
\begin{proof}
Put $T=\Q(E[3])$ and $F_0=\Q(\sqrt{-7})$.
The CM endomorphism action on $E[3]$ shows that $F_0\subset T$.
At the good prime three, which is inert in $F_0$, the local
CM reciprocity formula gives the full scalar inertia
$\mathbb F_9^\times=C_8$: the rank-one CM reflex norm is the
identity. Hence $\operatorname{Gal}(T/F_0)=C_8$.
Complex conjugation has order two and acts on this scalar
group by cubing, so $\operatorname{Gal}(T/\Q)=C_8\rtimes_3 C_2$.
These uses of good-prime reciprocity are the same
all-units formula as in \citet[Theorem~5.1, pp.~483--484]{Lombardo2017}.

The field $F(\sqrt[7]{3})$ has group $C_7\rtimes_3 C_6$:
the lift sending $\zeta_7$ to $\zeta_7^3$ and fixing the
chosen real seventh root has order six. Its intersection
with $T$ is precisely $F_0$, since both fields contain $F_0$
and their degrees are $42$ and $16$. Their compositum is
$K$ by Proposition~\ref{prop:lm-explicit-septic-field}.
Its Galois group is therefore the fiber product of these
two split groups over $C_2=\operatorname{Gal}(F_0/\Q)$.
Pairing the indicated order-six lift with complex conjugation
in the second factor gives an order-six lift in the fiber
product. It acts by cubing on both $C_7$ and $C_8$,
and therefore on $C_7\times C_8=C_{56}$. This proves the claim.
\end{proof}

\paragraph{An explicit polynomial specification.}
The three-division polynomial of $E$ is
\[
 f(X)=3X^4-3X^3-12X^2-12X-1.
\]
For $Z=2y+x$, put $g(X)=4X^3-3X^2-8X-4$.
An exact resultant computation gives
\[
 R(Z)=\operatorname{Res}_X(f(X),Z^2-g(X))
 =27Z^8-1323Z^6+6174Z^4-117649.
\]
The preceding subresultant is
\[
 (-36Z^4+1323Z^2-4116)X
       +9Z^4+588Z^2+1029.
\]
Its $X$-coefficient is coprime to $R(Z)$, so it recovers
$x$ from each root $z$, and then $y=(z-x)/2$.
It follows that $\mathbb Q(E[3])$ is the splitting
field of $R$.  Thus $K$ is explicitly the splitting
field of
\[
 (T^7-3)\bigl(27T^8-1323T^6+6174T^4-117649\bigr).
\]
The condition $\operatorname{Cl}(K)[3]=0$ will follow from the
degree-$56$ reduction and the principal-prime certificate below.
The defining polynomials alone do not establish class-group vanishing.
